\documentclass[11pt]{article}

\usepackage[T1]{fontenc}
\usepackage[utf8]{inputenc}
\usepackage[paper=a4paper,left=25mm,right=25mm,top=25mm,bottom=25mm]{geometry}
\usepackage[english, ngerman]{babel}
\usepackage{amsmath}
\usepackage{amssymb}
\usepackage[autostyle=true,german=quotes]{csquotes}
\usepackage{hyperref}
\usepackage{setspace}
\usepackage[expansion=false]{microtype}

\hypersetup{
    colorlinks,
    citecolor   = black,
    filecolor   = black,
    linkcolor   = black,
    urlcolor    = black,
    pdftitle    = {Expos\'e},
    pdfauthor   = {Erlind Sejdiu},
    pdfcreator  = {pdflatex},
    pdfproducer = {LaTeX with hyperref}
}
\setstretch{1.25}

\title{Expos\'e}
\author{Erlind Sejdiu}
\date{\today}

\begin{document}
\maketitle

\section{Motivation}

OpenTelemetry hat sich in den letzten Jahren als de facto Standard für Observability durchgesetzt. Es bietet eine umfassende Sammlung von APIs, SDKs und Tools zur Erzeugung, Sammlung und Analyse von Telemetriedaten wie Traces, Metrics und Logs.

Die einfache Integration von OpenTelemetry in Microservice-Architekturen und Cloud-nativen Anwendungen macht es zu einem unverzichtbaren Werkzeug für moderne Softwareentwicklung und -betrieb. In dieser Arbeit soll eine Analyse der bestehenden OpenTelemetry-Dokumentation und der verfügbaren Tools durchgeführt und die Integration in eine bestehende Microservice-Anwendung evaluiert werden. Ziel ist es, ein Verständnis für die Funktionsweise von OpenTelemetry und die Vorteile seiner Verwendung in der Praxis zu entwickeln.

Es existieren qualitiative Aussagen, wie OpenTelemtry eingesetzt werden kann, jedoch existieren keine quantitative Aussagen. Wie gut sich Opentelemetry in komplexen Systemarchitekturen, wie der im Unternehmen einsetzen lässt, ist unklar.

Deshalb soll in folgender Semniararbeit überprüft werden wie gut unterschiedliche Datenbanken unterschiedliche halten soll.

Das zuzügliche Logging verbraucht performanz als auch Speicherplatz. Eine Forschungsfrage ist die Analyse des Verbrauches und des Speicherplatzes.

Eine weitere Frage ist die Analyse von Wie und Wo unterschiedliche observabilities Möglichkeiten eingesetzt werden soll. Soll in jeder Funktion ein Traces entstehen. Ist dies im Debugging hilfreich, ist das Logging dann noch sinnvoll? Sollte man Event oder Attributes hinzufügen? Wie stark erfolgt

\section{Problemstellung}

\section{Stand der Technik}

\section{Forschungsfrage}

\section{Methodik}

\end{document}
